\section*{Введение}
\addcontentsline{toc}{section}{Введение}

В ходе курсовой работы был разработан программный комплекс для статического анализа растровых шрифтов в формате BDF. Программа читает BDF-файл, разбирает описание глифов, рассчитывает набор метрик для каждого символа, определяет потенциальные аномалии и формирует текстовый отчёт с предложениями замены на основе эталонного шрифта, заданного в конфигурации.

Формат BDF удобен для анализа в рамках функционального программирования, так как глифы представлены в текстовом виде: каждая строка пиксельной матрицы задаётся шестнадцатеричным числом. Это позволяет разобрать файл без внешних графических библиотек и далее работать с глифами как со списками строк.

Основная цель программы --- автоматизировать первичную проверку качества растрового шрифта. Анализатор оценивает читаемость, пропорциональность, плотность заполнения и различимость глифов. Если значение метрики выходит за допустимый диапазон, глиф помечается как аномальный. Для найденных аномалий программа ищет наиболее похожий глиф в эталонной базе и включает его в отчёт вместе с изображениями исходного и эталонного символа.

\newpage
